\section{Restricted tilting modules}
\label{sec:restricted}

Let $G=\mathrm{SL}_{2,k}$, let $V=k^2$ be its standard module, and set
\[
 u(t)=\begin{pmatrix}1&t\\0&1\end{pmatrix}.
\]
Let $U\simeq\mathbb G_a$ be the upper unitriangular subgroup,
parameterized by these matrices. Write $L(h)$ for the simple rational $G$-module of highest weight
$h\geq0$, and $\operatorname{St}_m=L(p^m-1)$ for the $m$th Steinberg
module; in particular $\operatorname{St}_0=k$.
The costandard and Weyl modules are
$\nabla(a)=\operatorname{Sym}^a(V^*)$ and $\Delta(a)=\nabla(a)^*$.
Here $V\simeq V^*$ by the invariant determinant pairing.
A module is \emph{tilting} if it has both a filtration by costandard
modules and a filtration by Weyl modules. Denote by $T(a)$ the
indecomposable tilting module of highest weight $a$.

\begin{theorem}\label{thm:restricted-tilting}
Let $r\geq1$, let $E_r=(\mathbb Z/p)^r$, and let
$\iota:E_r\hookrightarrow U(k)$ be any embedding. Write $R_\iota$
for restriction along $\iota$. For every finite-dimensional
$kE_r$-module $D$, there is a tilting $G$-module $T$ such that
\[
 R_\iota(\operatorname{St}_{r-1})\otimes D\simeq R_\iota(T).
\]
\end{theorem}

This proves the restricted-Steinberg assertion of
\cite[Conjecture~1]{CoulembierFlakeSteinberg}, for every embedding and
every prime. The proof has three steps. We lift $D$ to a vector bundle
on the plane with bounded highest weights at the origin. Tensoring
with $\operatorname{St}_{r-1}$ makes that origin fibre tilting. Finally,
an equivariant formal comparison identifies its restriction to $U$
with the fibre on the punctured plane. Write $R=R_\iota$ within this section.

\subsection{A vector bundle with bounded origin weights}
\label{subsec:geometric-lift}

For a finite-dimensional rational $U$-module $D$, its \emph{degree}
is the largest degree of an entry of its action matrix $\rho_D(u(t))$.
This number is independent of the chosen basis. The zero module will
always satisfy every degree bound below.

Keep the embedding $\iota:E_r\hookrightarrow U(k)$ fixed, and put $q=p^r$. We will replace an arbitrary
$kE_r$-module by a $G$-module at the origin of the standard plane.
The two modules need not be isomorphic on restriction; their comparison
will come after tensoring with a Steinberg module.

\begin{proposition}\label{prop:bounded-origin}
For every finite-dimensional $kE_r$-module $D$, there is a
$G$-equivariant vector bundle $\mathcal M$ on $\mathbb A^2_k$ such that
the stabilizer action on the fibre at $e_1=(1,0)$ restricts to $D$
along $\iota$. Every simple composition factor of the rational
$G$-module $M_0=\mathcal M|_0$ has highest weight less than $q$.
\end{proposition}

\begin{proof}
We first lift $D$ to a rational $U$-module of degree less than $q$.
Let $H\subset k$ be the image of $E_r$ under its additive embedding.
The vector space
\[
 C_q=k[t]_{<q}=\operatorname{span}_k\{1,t,\ldots,t^{q-1}\}
\]
is a subcoalgebra of $k[t]$, with comultiplication
$\Delta f(t,s)=f(t+s)$ and counit $f\mapsto f(0)$.
Evaluation induces a coalgebra map
\[
 \operatorname{ev}_H:C_q\longrightarrow k^H=\mathcal O(E_r).
\]
It is bijective: a polynomial of degree less than $q$ vanishing at
the $q$ distinct elements of $H$ is zero, and both spaces have
dimension $q$. Transporting the $\mathcal O(E_r)$-coaction of $D$
across this isomorphism and then including $C_q\subset k[t]$
gives a rational $U$-module $\widehat D$ of degree less than $q$.
Its restriction is $D$. This is a lift of individual representations;
we do not require the lifting operation to preserve tensor products.

The orbit map $G\to\mathbb A^2\setminus\{0\}$, $g\mapsto ge_1$,
is a right $U$-torsor. Indeed its fibres are right $U$-cosets,
and on the opens $x\ne0$ and $y\ne0$ it has the respective sections
\[
 \begin{pmatrix}x&0\\y&x^{-1}\end{pmatrix},\qquad
 \begin{pmatrix}x&-y^{-1}\\y&0\end{pmatrix}.
\]
Form the associated vector bundle
\[
 \mathcal E=G\times^U\widehat D,
 \qquad (g,d)\sim(gu,u^{-1}d),
\]
on this punctured plane, with $G$ acting by left multiplication.
The fibre identification $d\mapsto[1,d]$ at $e_1$ gives precisely
the prescribed $U$-action on $\widehat D$.

We next extend the bundle and its equivariance across the origin.
A coherent extension $\mathcal F$ of $\mathcal E$ exists on the
noetherian affine plane. Its double dual
$\mathcal M=\mathcal F^{**}$ is coherent and reflexive, and still
restricts to $\mathcal E$. On a regular surface a reflexive coherent
sheaf is locally free: its stalk has depth two at a closed point,
so the Auslander--Buchsbaum formula gives projective dimension zero;
the stalks of dimensions zero and one are free as well.
These standard extension statements can be found in
\cite[Tags 0G41, 0AY4, and 0B3N]{StacksProject}.

Write $a:G\times\mathbb A^2\to\mathbb A^2$ for the action and
$\pi_2$ for the projection. Over the punctured plane the bundle
action is an isomorphism $\pi_2^*\mathcal M\to a^*\mathcal M$.
Both sheaves are vector bundles on the regular variety
$G\times\mathbb A^2$. Their homomorphism sheaf is locally free,
and its sections extend uniquely across $G\times\{0\}$, which
has codimension two. This is the usual extension property for
reflexive sheaves \cite[Tag 0EBJ]{StacksProject}. Extend the isomorphism
and its inverse; their composites are the identity by uniqueness.
The identity and cocycle conditions hold on the punctured plane
and therefore, by the same uniqueness on $\mathbb A^2$ and
$G\times G\times\mathbb A^2$, hold everywhere. Thus
$\mathcal M$ is a $G$-equivariant vector bundle.

It remains to bound the highest weights of $M_0$. The line
$\ell=ke_1$ is fixed pointwise by $U$. For $a\in k^\times$ put
$h_a=\operatorname{diag}(a,a^{-1})$. The bundle action of $h_a$
identifies the fibres at $e_1$ and $ae_1$, and
\[
 h_a^{-1}u(t)h_a=u(a^{-2}t).
\]
Consequently the $U$-action on every nonzero fibre of
$\mathcal M|_\ell$ has degree less than $q$.

To pass this bound to the origin, let
$A=\Gamma(\ell,\mathcal End(\mathcal M|_\ell))$.
Since $U$ acts trivially on the base $\ell$, its bundle action is
an element of $A\otimes_k k[t]$, say
$\sum_{j=0}^{N} b_jt^j$. For $j\geq q$, the section $b_j$
vanishes on every nonzero fibre. The field $k$ is infinite and
$\mathcal End(\mathcal M|_\ell)$ is locally free, so $b_j=0$.
Specialization at zero therefore gives an action on $M_0$ of
degree less than $q$. This bound passes to submodules and
quotients by writing the action in a basis adapted to a submodule.

Finally, if $L(h)$ is the simple $G$-module of highest weight $h$,
its $U$-action has degree exactly $h$. To see this, write
$h=\sum_i h_ip^i$ with $0\leq h_i<p$. Steinberg's tensor product
theorem, in the form recalled in \cite[Lemma 1.1 and the proof of
Theorem 2.1]{DotyHenke2005}, gives
\[
 L(h)\cong\bigotimes_i
 \bigl(\operatorname{Sym}^{h_i}(k^2)\bigr)^{(i)}.
\]
Here $(i)$ denotes Frobenius twist. The action on the $i$th factor
has degree $h_ip^i$. On the tensor product, the coefficient from
$\bigotimes_i e_2^{h_i}$ to $\bigotimes_i e_1^{h_i}$ is $t^h$,
so the upper bound $h$ is attained. Hence every simple factor
$L(h)$ of $M_0$ satisfies $h<q$, as required. The argument applies
unchanged when $p=2$.
\end{proof}


\subsection{The origin fibre after tensoring with a Steinberg module}
\label{subsec:bounded-tilting}

We next convert the highest-weight bound on the origin fibre into a tilting
module. We use the theory of tilting modules developed by Donkin
\cite{Donkin1993}. Tilting modules are closed under extensions and tensor products.
For $0\leq a<p$ one has $L(a)=T(a)$, and Donkin's formula is
\begin{equation}\label{eq:donkin-sl2}
 T(a)\otimes T(b)^{(1)}\simeq T(a+pb)
 \qquad(p-1\leq a\leq2p-2,\quad b\geq0),
\end{equation}
where $(1)$ means pullback along the Frobenius morphism of $G$.
The identity $L(a)=T(a)$ in the stated range and Formula~\eqref{eq:donkin-sl2}
hold in every positive characteristic, including $2$;
see \cite[Proposition~3.1]{BEO2023}.

The following is the extension-closed form of
\cite[Lemma~3.3]{BEO2023}, also used in the proof of
\cite[Lemma~3.1.6]{CoulembierFlakeSteinberg}.  We recall its short proof
to make the highest-weight bound explicit.

\begin{lemma}\label{lem:bounded-tilting}
Let $k$ be algebraically closed of characteristic $p>0$, let $m\geq0$,
and let $M$ be a finite-dimensional rational $SL_2(k)$-module.
If every simple composition factor of $M$ is $L(h)$ with
$0\leq h<p^{m+1}$, then $\operatorname{St}_m\otimes M$ is tilting.
Here $\operatorname{St}_0=k$.
\end{lemma}

\begin{proof}
Tensoring a composition series of $M$ by $\operatorname{St}_m$ and splicing
Weyl and costandard filtrations reduces the assertion to $M=L(h)$.
We induct on $m$.  When $m=0$, the inequality $h<p$ gives
$\operatorname{St}_0\otimes L(h)=T(h)$.
For $m>0$, write $h=d+ph'$ with $0\leq d<p$ and $h'<p^m$.
Steinberg's tensor product theorem gives
\[
 \operatorname{St}_m\otimes L(h)
 \simeq
 L(d)\otimes\operatorname{St}_1\otimes
 \bigl(\operatorname{St}_{m-1}\otimes L(h')\bigr)^{(1)}.
\]
By induction, the module in parentheses is a direct sum of tilting
modules $T(b)$.  Formula~\eqref{eq:donkin-sl2}, applied with $a=p-1$,
shows that tensoring its Frobenius twist with $\operatorname{St}_1$
is tilting.  Tensoring once more with $L(d)=T(d)$ preserves that property.
\end{proof}

Apply Lemma~\ref{lem:bounded-tilting} to the origin fibre $M_0$ with
$m=r-1$.  Its established bound $h<p^r$ shows that
\[
 T:=\operatorname{St}_{r-1}\otimes M_0
\]
is tilting.  This includes $r=1$, when the Steinberg factor is trivial.
It remains to compare this origin fibre with the given module on the
punctured plane.

\subsection{Comparing the origin and the unipotent fibre}
\label{subsec:formal-comparison}

The bundle $\mathcal N=\operatorname{St}_{r-1}\otimes_k\mathcal M$
has the tilting origin fibre $T$ constructed above.  We now recover its
representation on the punctured orbit from that fibre.  The following
lemma applies to any equivariant bundle with tilting origin fibre:
we first trivialize it equivariantly on successive infinitesimal
neighbourhoods of $0$, then compare the generic point of the $U$-fixed
line with $e_1$.

\begin{lemma}[Comparison with a tilting origin fibre]
\label{lem:formal-comparison}
Let $k$ be an algebraically closed field of arbitrary positive
characteristic.  Let $\mathcal N$ be a $G$-equivariant vector bundle on $V$,
and suppose that its origin fibre $T=\mathcal N_0$ is a tilting
$G$-module.  Then
\[
 \mathcal N_{e_1}\cong T|_U
\]
as rational $U$-modules over $k$.
\end{lemma}

\begin{proof}
We first construct compatible equivariant trivializations on every
infinitesimal neighbourhood of $0$.  Let $A=k[x,y]$, let
$\mathfrak m=(x,y)$, and let $N=\Gamma(V,\mathcal N)$, a possibly
infinite-dimensional rational $G$-module.  Its quotients
$N/\mathfrak m^jN$ are finite dimensional.  Local freeness gives
canonical $G$-isomorphisms
\begin{equation}
 \label{eq:jet-kernel}
 \mathfrak m^jN/\mathfrak m^{j+1}N
 \cong (\mathfrak m^j/\mathfrak m^{j+1})\otimes_k T
 \cong \operatorname{Sym}^j(V^*)\otimes_k T
 \qquad(j\geq 0).
\end{equation}

Recall that a good filtration has costandard successive quotients, and a
Weyl filtration has standard successive quotients.  For $G=\mathrm{SL}_2$,
the costandard module of highest weight $a$ is
$\nabla(a)=\operatorname{Sym}^a(V^*)$.  Tensor products of two such modules
have good filtrations.  Indeed, for $a,b\geq1$, multiplication of polynomials
gives the exact sequence
\begin{equation}
 \label{eq:symmetric-power-filtration}
 0\longrightarrow
 \operatorname{Sym}^{a-1}(V^*)\otimes\operatorname{Sym}^{b-1}(V^*)
 \xrightarrow{\;d\;}
 \operatorname{Sym}^{a}(V^*)\otimes\operatorname{Sym}^{b}(V^*)
 \longrightarrow \operatorname{Sym}^{a+b}(V^*)
 \longrightarrow0.
\end{equation}
Here $d=x_1y_2-y_1x_2$ is the determinant in two copies of $V^*$.
Multiplication by $d$ is injective, its image lies in the kernel, and
dimensions prove exactness.  Induction on $\min(a,b)$, starting with a
trivial factor, gives the assertion.  Refining filtrations then proves the
same tensor-product assertion for arbitrary modules with good
filtrations.  This argument is unchanged in characteristic two.

The module $T$ has both a Weyl filtration and a good filtration.
Consequently, the target of the next Ext group has a good filtration,
and the standard--costandard Ext vanishing theorem of
Cline--Parshall--Scott--van der Kallen gives
\begin{equation}
 \label{eq:formal-ext-vanishing}
 \operatorname{Ext}^1_G
 \bigl(T,\operatorname{Sym}^j(V^*)\otimes T\bigr)=0
 \qquad(j\geq0).
\end{equation}
We use Ext in the category of rational $G$-modules; the vanishing follows
from $\operatorname{Ext}^1_G(\Delta(a),\nabla(b))=0$ by induction on the two
filtrations; see \cite[Chapter~1, Proposition~2.6, and Chapter~4,
Proposition~1.1]{Riche}.

Starting with the identity $s_1:T\to N/\mathfrak mN$, lift successively to
$G$-maps
\[
 s_j:T\longrightarrow N/\mathfrak m^jN\qquad(j\geq1)
\]
such that $s_{j+1}$ reduces to $s_j$.  Pull back the short exact sequence
\[
 0\longrightarrow \operatorname{Sym}^j(V^*)\otimes_k T
 \longrightarrow N/\mathfrak m^{j+1}N
 \longrightarrow N/\mathfrak m^jN\longrightarrow0
\]
along $s_j$.  By \eqref{eq:formal-ext-vanishing} the resulting extension
of $T$ splits equivariantly; a splitting gives the lift $s_{j+1}$.
Extending $s_j$ by scalars gives an equivariant
$A/\mathfrak m^j$-linear map
\[
 (A/\mathfrak m^j)\otimes_k T\longrightarrow N/\mathfrak m^jN.
\]
Both modules are free of rank $\dim_kT$ over the local Artinian ring
$A/\mathfrak m^j$, and the map reduces to the identity on the residue
field.  Nakayama's Lemma therefore makes it an isomorphism.  Taking the
inverse limit produces an isomorphism
\begin{equation}
 \label{eq:formal-trivialization}
 k[[x,y]]\otimes_kT\xrightarrow{\sim}\widehat N
\end{equation}
compatible with the $G$-action at every finite order.

The line $y=0$ is fixed scheme-theoretically by $U$.  Restrict
\eqref{eq:formal-trivialization} to this line and then set $K=k((x))$.
This gives a rational $U_K$-isomorphism
\[
 \varphi:T_K\xrightarrow{\sim}\mathcal N_{xe_1}.
\]
To justify the group-scheme assertion explicitly, choose a polynomial
frame for $\mathcal N$ on the line; this is possible because a finite
projective $k[x]$-module is free.  The action is a matrix
$C(x,t)\in\mathrm{GL}_{\dim T}(k[x,t])$, and the action on $T$ is a
matrix $B(t)\in\mathrm{GL}_{\dim T}(k[t])$.  If $H(x)$ is the matrix of
the restricted formal isomorphism, equivariance at every finite order
says
\[
 C(x,t)H(x)=H(x)B(t)\pmod{x^j}
 \qquad\text{for every }j\geq1.
\]
These matrices have uniformly bounded degree in $t$, so coefficientwise
$x$-adic separation gives the equality in $k[[x]][t]$, and hence over
$K[t]$.

We must compare the two fibres without changing the parameter of $U$.
Put $h=\operatorname{diag}(x,x^{-1})\in G(K)$, and let
$b_h:\mathcal N_{e_1,K}\to\mathcal N_{xe_1}$ be bundle transport.
Write $\rho_T$ for the $G$-action on $T$ and $\rho_{e_1},\rho_x$ for
the two fibre actions of $U$.  Since
$h u(t)h^{-1}=u(x^2t)$, equivariance gives
\[
 b_h\rho_{e_1}(u(t))=\rho_x(u(x^2t))b_h,
 \qquad
 \rho_T(h)\rho_T(u(t))=\rho_T(u(x^2t))\rho_T(h).
\]
It follows that
\begin{equation}
 \label{eq:cancel-torus-conjugation}
 f=b_h^{-1}\varphi\rho_T(h):T_K\xrightarrow{\sim}\mathcal N_{e_1,K}
\end{equation}
intertwines the original $U$-actions.  Thus restriction of $f$ to any
fixed finite subgroup of $U(k)$ uses that subgroup's original embedding.

Finally, this isomorphism descends to $k$.  In fixed $k$-bases, the
intertwiner space
\[
 L=\operatorname{Hom}_U(T,\mathcal N_{e_1})
\]
is the kernel of finitely many linear equations over $k$: compare the
coefficients of $t$ in the two polynomial action matrices.  Hence its
scalar extension to $K$ is the intertwiner space over $K$.
Equation~\eqref{eq:cancel-torus-conjugation} shows that the determinant
polynomial on $L$ is nonzero.  Since $k$ is infinite, this polynomial is
nonzero at some point of $L(k)$, giving the required isomorphism.  At no
point have we asserted that the formal trivialization algebraizes on
the plane.
\end{proof}

\begin{proof}[Proof of Theorem~\ref{thm:restricted-tilting}]
Given $D$, choose $\mathcal M$ as in Proposition~\ref{prop:bounded-origin}
and put $M_0=\mathcal M|_0$. Lemma~\ref{lem:bounded-tilting} makes
$\operatorname{St}_{r-1}\otimes M_0$ tilting. Apply Lemma~\ref{lem:formal-comparison} to
$\mathcal N=\operatorname{St}_{r-1}\otimes_k\mathcal M$.  Its origin
fibre is the tilting module
$T=\operatorname{St}_{r-1}\otimes_kM_0$.  By
Proposition~\ref{prop:bounded-origin}, the fibre of $\mathcal N$ at $e_1$
restricts to $R(\operatorname{St}_{r-1})\otimes D$ along the fixed
embedding $E_r\hookrightarrow U(k)$.  The comparison therefore gives
\[
 R(T)\cong R(\operatorname{St}_{r-1})\otimes D,
\]
as required.  This comparison also applies when $r=1$, with
$\operatorname{St}_0=k$.
\end{proof}
